\section{Симуляция}

\tab Одним из самых популярных методов функциональной верификации является симуляция. Основная идея этого метода заключается в создании особого окружения (тестбенча) вокруг проверяемого устройства (DUT). Это окружение должно генерировать векторы входных данных, после чего подавать эти векторы на вход устройства. Одновременно с этим окружение производит независимые вычисления, преобразовывая входные векторы в соответствии со спецификацией. Далее окружение собирает выходной вектор с устройства и сравнивает его с предсказанным вектором.

// Вставить картинку

\subsection{Направленное тестирование}

\tab Самым примитивным подходом динамической верификации является направленное тестирование. Используя его, инженер на основе спецификации составляет список тестом, каждый из которых фокусируется на одном отдельном свойстве устройства. На основе этого плана практически вручную прописываются входные векторы которые тестируют заданные свойства. После успешного прохождения теста, он вычеркивается из полного списка, и инженер переходит к следующему пункту. 

Такая инкрементивность гарантирует почти мгновенные результаты. Тем не менее, она же и делает этот подход крайне немасштабируемым: если сложность устройства удваивается, то и время на его верификацию увеличивается вдвое. В связи с этим ограничением направленное тестирование как основной способ верификации сейчас практически не используется. 

\subsection{Ограниченное случайное тестирование}

\tab В настоящее время самый популярный подход --- ограниченное случайное тестирование. В случае использования данного подхода входные векторы не прописываются инженером самостоятельно, а генерируются случайно, хоть и в рамках заданных ограничений и весовых коэффициентов. Основной инженерной сложностью становится разработка автоматизированного способа предсказывать результат, который, как правило, представляет собой референсную модель устройства.

Изначальная задержка в получении результата ввиду построения референсной модели полностью компенсируется в дальнейшем: единый общий тестбенч для всех проводимых тестов находит ошибки гораздо быстрее, чем множество тестбенчей в направленном тестровании. 

В связи с высокой масштабируемостью, удобством использования и быстрым получением результатов, ограниченное случайное тестирование является домирнирующим способом верификации уже многие годы. Ввиду этого, было предложено и разработано множество методик, например, OVM, VVM, UVM и др. В дальнейшем в работе в качестве симуляции как метода верификации будет рассматриваться только UVM ввиду своей популярности. 

\subsubsection{UVM}


Ключевыми компонентами являются драйвер, монитор, секвенсеры, агенты, скорборд. 

\subsubsection{Покрытие}

 Функциональное покрытие разрабатывает сам верификатор. Он устанавливает диапазоны входных и выходных сигналов, а также их комбинации, которые должны быть проверены в тестах. Кодовое покрытие собирается автоматически инструментами верификации. 

В свою очередь, оно подразделяется на 4 подвида: 
\begin{itemize}
    \item Block.
    \item Statement. 
    \item Expression.
    \item Toggle. 
\end{itemize}

\subsection{Трассы}

Как правило, при разработке сложных микроархитектур вначале пишется не RTL-код, а потактовый симулятор разрабатываемого устройства на языке более высокого уровня(например, C++). Основное назначение такого симулятора --- составить референсную модель. С помощью этой рефенсной модели можно проанализировать происзводительность на сложных и длинных тестах, которые затрачивали бы нереаличтичные промежутки времени на потактовом симуляторе RTL. 

Вместе с тем, в процессе анализа формируются векторы входных с выходных данных, которые проходят через эту модель. Их будем называть трассами. Для проверки соответствия разработанного ртл модуля можем использовать эти трассы как средство верификации. 

Практически таким же образом можем использовать чекпоинты --- снимки состояния системы в определённые моменты времени. 



\subsection{Ключевые особенности}

Симуляция имеет несколько ключевых особенностей. Во-первых, т.к. мы выборочно “обстреливаем” блок случайными данными, то невозможно с уверенностью сказать, что блок работает корректно, на основании одного лишь теста. Поэтому, как правило, после разработки окружения и сбора покрытия, т.е. После того, как мы удостоверились, что тесты покрывают все возможные случаи поведения блока, начинается этап регрессии. Регрессия – это регулярный запуск тестов (как правило, каждую ночь), с целью удостовериться в корректности блока. 

Другой ключевой особенностью является отсутствие заметной зависимости времени симуляции одного такта от сложности блока. Таким образов, сложная вычислительная логика, последовательная логика, широкие шины или память не представляют особых затруднений для основанных на симуляции методов. Это привносит важное качество универсальности --- метод можно применять ко всем уровням абстракции: от мелких компонент до полноценной СнК. 

В-третьих, верификатор должен прилагать усилия, чтобы данные были достаточно разнообразными. 

Невозможно верифицировать блок раньше напсания, зависимость от правильной работы протоколов передачи

